\begin{figure*}[t]
\centering
\scriptsize
\setlength{\fboxsep}{1em}
\setlength{\fboxrule}{1pt}
\fcolorbox{black}{white}{
\parbox{0.955\textwidth}{
\begin{alltt}


You will be provided instructions for completing a task followed by a task to complete.

Instructions:

1. Understand the structure of the prompt: The prompt will present an article with a title, introduction, background information, main argument, supporting evidence, counterargument, and conclusion.

2. Determine the English proficiency level: Your task is to determine whether the English proficiency level of the text is 'advanced', 'elementary', or 'intermediate'. 

3. Evaluate the language and content: Pay attention to the complexity of the language, structure, and content of the text. 

- 'Elementary': Look for simple language and basic phrases. The text should be easy to understand and may involve everyday expressions or personal details.

- 'Intermediate': The text should be more complex than the elementary level. Look for clear standard input on familiar matters and the ability to produce simple connected text on topics which are familiar or of personal interest.

- 'Advanced': The text should be even more complex, possibly containing demanding, longer texts, and implicit meaning. Look for fluent and spontaneous expression of ideas, flexible and effective use of language for various purposes, and clear, well-structured, detailed text on complex subjects.

4. Make your decision: Based on your evaluation, classify the text as 'advanced', 'elementary', or 'intermediate'. 
Remember, these levels are often determined based on the Common European Framework of Reference for Languages (CEFR), an international standard for describing language ability.

Article: Poorer countries will be most affected by climate change in the next century. 

Sea levels will rise, there will be stronger cyclones, warmer days and nights, more rainfall, and larger and longer heatwaves, says a new report. 

The last big United Nations (UN) report, in 2007, said there would be temperature rises of 6°C or more by the end of the century. Scientists now think this will not happen, but average land and sea temperatures will probably continue rising during this century, possibly becoming 4 °C hotter than now. That rise would ruin crops and make life in many cities too hot. 

As temperatures rise and oceans become warmer, there will be big changes in annual rainfall in tropical and subtropical regions, says the Intergovernmental Panel on Climate Change (IPCC) report, released in Stockholm and published online in September 2013. 

East Africa can expect more short rainfalls and west Africa should expect heavier monsoons. Burma, Bangladesh and India can expect stronger cyclones; elsewhere in southern Asia, there will probably be heavier summer rainfalls. Indonesia may receive less rainfall between July and October, but the coastal regions around the south China Sea and Gulf of Thailand can expect more rainfall when cyclones hit the land. 'Rainfall patterns will change. Northern countries, for example in Europe or North America, will probably receive more rainfall, but many subtropical, dry regions will likely get less rain,' said the report. 

The report also said that the monsoon season will probably start earlier and last longer. 

Scientists in developing countries are happy with the report. “The IPCC says that climate change is real and happening much more strongly than before. We are already seeing the effects of climate change in Bangladesh and across south Asia. It’s not news to us. 

Most developing countries are experiencing climate change now. They do not need the IPCC to tell them that the weather is changing,” said Saleemul Huq, director of the International Centre for Climate Change and Development. 

Scientists have also lowered their predictions for sea-level rises. Sea levels will probably rise an average of 40 –62 cm by 2100. But many millions of people living in the developing world’s great cities, including Lagos and Calcutta, are in danger. 

Weather disasters are also more likely in a warmer world, the report says. The number of tropical cyclones will probably not change, but they may become more intense, with stronger winds and heavier rainfall. 

Life in many developing country cities could become very difficult, especially because city temperatures are already higher than those in the countryside. Much higher temperatures could reduce the length of the growing period in some parts of Africa by up to 20%

The charity Oxfam said that world hunger would get worse because climate changes hurt crop production. They said the number of people at risk of hunger might increase by 10%

“The changing climate is already hurting the fight against hunger, and it looks like it will get worse,” said Oxfam. “A hot world is a hungry world”.

Follow the instructions to generate an explanation that reasons towards the correct answer to the task above. End the explanation with the correct answer. [/INST] Explanation: 

The article provided discusses the impact of climate change on poorer countries in the next century. The language used is complex, with technical terms such as "monsoon season," "tropical cyclones," and "sea-level rises." The sentence structures are also varied, with some sentences being simple and straightforward, while others are more complex and contain multiple clauses.

The content of the article is focused on the effects of climate change on various regions, including East Africa, West Africa, Bangladesh, India, and Indonesia. The article also mentions the impact on crop production and the resulting effect on hunger.

Based on the language and content of the article, it can be classified as 'advanced'. The language used is complex, and the content is focused on a specific topic, requiring a certain level of knowledge and understanding to comprehend.

Therefore, the correct answer is:

Advanced

Therefore, the correct label among \{'advanced', 'elementary', 'intermediate'\} (just the label) to the original task is below.

Label: Advanced
\end{alltt}
}}
\caption{A complete running example of RAFT (One Stop English).}
\label{fig:raft_subset=one_stop_english_example}
\end{figure*}
